\subsection{Resumen integrado de esfuerzo y capacidad}
Synaptix selecciona 49.396 HH en 2181 paquetes y 122.070 HH de cobertura propia: 171.466 HH profesionales. Los 1.959 paquetes de origen y sus 40.744 HH se conservan; 216 paquetes de conciliación funcional añaden 8.508 HH. Se añaden seis paquetes mensuales de mentoría, 144 HH, con dos puestos parciales propios e independientes de cobertura. Los padres no reciben horas. Los códigos enlazan resultado, responsable, perfiles, actividades, red y recepción; las reservas representan capacidad futura y no trabajos ejecutados.
La Tabla~\ref{tab:sd7-esfuerzo-integral} presenta esfuerzo de entregables y cobertura por fase.
\begin{table}[H]
\centering
\begin{tabularx}{\linewidth}{Y R{25mm}}\hline \EncabezadoTabla{Concepto} & \EncabezadoTabla{HH}\\\hline
Implementación: entregables e ingeniería & 30.872\\\hline
Cobertura inicial y marchas blancas (13--20) & 22.692\\\hline
Operación: entregables e ingeniería & 18.524\\\hline
Operación: NOC/SOC y mesa (21--56) & 99.378\\\hline
Total profesional programado & 171.466\\\hline
\end{tabularx}
\caption{Esfuerzo de entregables y cobertura por fase}
\label{tab:sd7-esfuerzo-integral}
\FuenteTabla{Elaboración propia: diccionario y asignaciones integradas de T-14/T-15.}
\end{table}\par\smallskip
La suma distingue 49.396 HH de entregables e ingeniería de 122.070 HH de cobertura; ambas partidas concilian las 171.466 HH sin duplicar turnos. Ingeniería dispone de 7,2 HH productivas por día y un máximo de 50 HH de paquetes por persona y quincena. Los turnos y el control natural final disponen de ejecutores propios distintos, con máximo de 144 HH al mes, 40 HH por siete días y doce horas de descanso. El máximo mensual de dotación comprometida es 106 puestos, incluidos los ocho titulares; el máximo de refuerzos de ingeniería es 98. Es el resultado de esta asignación, sin atribuir contratación ni optimización global. Synaptix proveerá estaciones, accesos, entornos, incorporación y suplencias para cada puesto antes de su intervención; los ocho equipos de los titulares no cubren la dotación adicional.
La habilitación común inicial se imputa a E1 y el gobierno posterior al mes 12 al desarrollo/cierre E2. La estabilización, aceptación final y cobertura de E1 en producción se mantienen como partidas adicionales identificadas. Los máximos por fase no se suman como personas diferentes. La valoración profesional se desarrolla exclusivamente en la Oferta Económica.
